\section{问题重述与总体分析}
\subsection{研究对象与任务}
红外反射光谱中的条纹记录了光在外延层内往返传播时累积的相位。相位随波数的变化速度同时取决于层厚、材料色散和折射后的传播方向。因而，条纹周期能够约束厚度，但周期与几何厚度之间还需要一条经过说明的折射率关系。对掺杂外延材料，吸收和界面反射也会改变条纹形状，不能把每一个反射率极值都等同于固定相位条件。

本文处理两种材料、每种材料两个入射角的四条实测光谱。附件1和附件2来自同一碳化硅晶圆，入射角为$10^\circ$和$15^\circ$；附件3和附件4来自同一硅晶圆，入射角相同。输入为波数与反射率，不包含独立的层厚测量、掺杂浓度、晶型、晶向、温度、光斑位置及仪器响应函数。因此，研究目标是给出\emph{在明确光学假设下可以复现的厚度估计}，同时识别数据尚不能确定的部分。

问题一要求在只保留一次界面反射和透射的条件下建立厚度模型，需要区分几何光程、同一出射波前上的相位差及反射相移。问题二要求将模型用于碳化硅光谱，并说明预处理、求解方法和可靠性。问题三要求把界面上的多次往返传播纳入模型，讨论多光束干涉形成与可观测的条件，再分析硅光谱，最后检查碳化硅结果是否需要相应修正。

\subsection{从条纹定位到模型比较}
我们把求解划分为三个层次。首先检查原始数据，并通过条纹频率确定光学厚度的大致范围，避免非线性拟合落入错误条纹级次。其次引入与波段相符的折射率关系，在每个材料内部比较两个角度的估计，并考察同一厚度约束能否解释两条曲线。最后在相同拟合窗口和相同数据口径下比较两光束与多光束描述，以预测误差、残差结构和厚度变化共同判断修正的必要性。

\paperfigure{figures/workflow.png}{由材料核对到条件厚度估计的分析流程}{0.78}{fig:workflow}
\figref{fig:workflow}中的可靠性检查同时作用于数据和模型。改善曲线拟合不自动代表厚度更准确：增加参数可以减小残差，却可能使折射率、吸收和厚度之间的补偿更严重。本文保留简单基准，使模型复杂度增加所带来的收益可以检查。

\subsection{研究假设与适用边界}
\begin{enumerate}
  \item 测量光斑内的外延层近似平面平行、横向均匀，忽略宏观曲率；两个角度来自同一晶圆，但不把附件说明进一步解释为已经证实光斑严格重合。
  \item 入射环境折射率近似为1，题给角度相对于界面法线定义；角度误差另作灵敏度分析，不在主拟合中随意调整角度。
  \item 相邻采样点对应同一测量配置，波数轴单调且单位为$\mathrm{cm}^{-1}$。保留反射率的百分数口径，拟合误差统一报告为百分点。
  \item 在选定的透明或弱吸收波段内，可用文献色散关系描述实折射率；文献关系不等同于对本晶圆的独立折射率测量，其偏差应计入模型不确定性。
  \item 缓慢变化的仪器增益、背景和干涉包络可用低阶函数描述。其变化尺度应大于条纹周期，不以高阶背景函数吸收条纹。
  \item 多光束模型将可相干的往返反射进行振幅叠加；散射、厚度分布及仪器卷积可能降低可见度，在缺少独立测量时不分别反演所有机制。
  \item 残差可包含沿波数相关的系统成分。因此，不把全部密集采样点视为独立重复实验，也不将条件统计区间解释为绝对测量准确度。
\end{enumerate}
这些假设决定了结果的解释范围。若材料晶型、自由载流子贡献或测点厚度差异未被正确描述，即使求解器收敛，所得几何厚度仍可能发生系统偏移。

\subsection{主要符号}
\begin{table}[htbp]\centering
\caption{主要符号及单位}\label{tab:symbols}
\begin{tabularx}{\textwidth}{cLc}
\toprule 符号 & 含义 & 单位\\\midrule
$\wn$ & 真空波数 & $\mathrm{cm}^{-1}$\\
$\lambda$ & 真空波长 & $\mu\mathrm{m}$\\
$d$ & 外延层几何厚度 & $\mu\mathrm{m}$\\
$\theta,\beta$ & 外部入射角和层内折射角 & 度或弧度，按公式说明\\
$\widetilde n=n+\ii\kappa$ & 复折射率及消光系数 & 1\\
$\delta$ & 单程传播相位 & rad\\
$\Phi$ & 两束光的往返相位差 & rad\\
$r_{ij},t_{ij}$ & 界面振幅反射和透射系数 & 1\\
$R,y$ & 模型反射率和实测反射率 & 1或\%，按上下文说明\\
$X_\theta(\sigma)$ & 色散修正后的相位坐标 & $\mathrm{cm}^{-1}$\\
$\boldsymbol\eta$ & 背景、包络及相移等辅助参数 & 依具体参数定义\\
$\widehat d$ & 给定模型下估计的厚度 & $\mu\mathrm{m}$\\
\bottomrule
\end{tabularx}\end{table}
\tabref{tab:symbols}只列全局符号。数值计算使用微米表示厚度时，传播相位中的$10^{-4}$用于将微米换成厘米；这个转换在代码与公式中保持一致。
